This revision uses four evidence classes.
v6 adds the last one.
\begin{description}[style=nextline,leftmargin=0pt]
\item[Locally proved.]
The normalized scalar apex coefficient; the singular-channel heat
shift; the new \(r^2\)-weighted coefficient and its shift; a
conditional graph-reduction theorem; the longitudinal UV
normalization; algebraic identities between the printed polynomials;
and the stated access/flux distinctions.

\item[Executed and artifact-verified in v6.]
Statements carrying the tag \textsc{[exec]}.
Each was produced by a program in \loc{code/} and logged in
\loc{logs/v4/} or \loc{logs/v6/}, and was rerun or checked in this
revision.
Its full SHA-256 digest is listed in Appendix~\ref{app:artifacts}.
The vertex inputs are read-only copies of the frozen XVI-P/Q
files; their digests equal the originals.
Exact statements are exact rational or symbolic identities in
\(\alpha\) or \(\nu\).
Numerical statements give their precision.
An executed result is a reproducible computation, not an
independent review.

\item[Supplied computational record.]
v4 reports that v6 did not rerun, together with the historical
registry intake.
These are kept as reports.

\item[Conditional graph identification.]
Equality of a printed polynomial with a particular frozen graph
integral still needs the graph-specific hypotheses of
Assumption~\ref{ass:reduction}.
v6 discharges several of these for the constant-profile record;
the ledger in Section~\ref{sec:ledger} lists them.
No proof is reconstructed from a routine name.
\end{description}

\scopebox{Submission and proof boundary (v6)}{%
v6 now archives the vertex arrays, routing, prefactors, and logs
that v5 found missing.
The constant-profile catalogue is therefore an executed result on
the archived frozen inputs, and no longer just a supplied report.
It is still not an independently refereed computation.
The spectral reduction of each graph remains conditional on
Assumption~\ref{ass:reduction}.
Only (R1) and parts of (R3) are discharged here
(Sections~\ref{sec:catalogue}, \ref{sec:HI}).
(R2), (R4), and (R5) are still the stated hypotheses.}

The coefficient values printed in v4 and v5 are unchanged.
v6 adds entries but alters none.

\subsection{What v6 adds}

\begin{enumerate}[label=\textup{(\alph*)}]
\item \emph{\(hh\) Slavnov--Taylor row}
(Section~\ref{sec:ST}).
Derived conventions, an \emph{a priori} constant \(\lambda\), and
exact closure on CR and F for the constant profile.
The v4 failure is diagnosed, and UNIT closes.
\item \emph{Higher apex coefficient}
(Section~\ref{sec:weighted}).
The \(r^2\)-weighted coefficient \(\mathcal D_b\), with checks.
The quadratic-profile \(\phi\phi\) value is partial: the D05a part
is \(0\) and the D04 part is open.
\item \emph{\((\mathrm H_{I,\mathrm{D01}})\)}
(Section~\ref{sec:HI}).
The exact \(L^2\) obstruction, and an exact vanishing Mellin value.
\item \emph{Evanescent coefficients}
(Section~\ref{sec:missing}).
A pole-level statement only; the coefficients themselves are open.
\item \emph{\(hh\) map}
(Section~\ref{sec:structures}).
Six invariant structures, the complete D02/D12 arrays, and the
coefficient vectors.
The 18-dimensional engine map is open.
\item \emph{\(\Span\) certificate}
(Section~\ref{sec:missing}).
Open, with the precise reason.
\item \emph{Response table}
(Section~\ref{sec:response}).
Every v5 issue, answered point by point.
\end{enumerate}
